%% chapter4.tex — Experimental Evaluation (RA-TDMAs+ Layer-3 routing)
%% Versão completa e detalhada, alinhada com os dados medidos no testbed de 3 nós.
%% As figuras já carregadas no Overleaf:
%%   RTT_over_the_direct_link.png   (\ref{fig:rtt-direct})
%%   RTT_over_the_relay_path.png    (\ref{fig:rtt-relay})
%%   timing_tdma.png / IA histograms (\ref{fig:ia-direct}, \ref{fig:ia-relay})
%%   slot occupancy / beacons       (\ref{fig:slots})

\chapter{Experimental Evaluation}
\label{chap:evaluation}

\section{Introduction}
\label{sec:eval-intro}

This chapter presents the experimental evaluation of the Layer~3 routing
system described in Chapter~\ref{chap:implementation}. The evaluation was
carried out on physical hardware in a three-node testbed that reproduces the
target operational scenario: a mobile robot transmitting a live video stream
to a stationary base station, with an intermediate relay node available to
maintain connectivity after a direct-link failure.

The evaluation pursues two objectives. The first is to characterise the
end-to-end performance of the system --- control-channel latency, video
throughput, and recovery time after a link failure. The second, and central
to this dissertation, is to \emph{prove that the nodes respect their TDMA
slots} when transporting both periodic synchronisation beacons and a
continuous video stream, and that this property is preserved when traffic is
forwarded through the relay.

The chapter is organised as follows.
Section~\ref{sec:setup} describes the hardware, the per-node network and
software configuration, the video pipeline, and the link-failure methodology.
Section~\ref{sec:metrics} defines the performance metrics.
Section~\ref{sec:experiments} describes the five experiments performed and
exactly how the data was collected.
Section~\ref{sec:results} presents and analyses the results.
Section~\ref{sec:comparison} compares the system with the Layer~2 approach of
Morais~\cite{morais2024synchronization}.
Section~\ref{sec:eval-summary} summarises the findings.


%% =======================================================================
\section{Experimental Setup}
\label{sec:setup}

\subsection{Testbed Overview}
\label{subsec:overview}

The testbed comprises three nodes, each with a fixed role. The network is an
IEEE~802.11 ad-hoc (IBSS) cell on channel~6. Two addressing planes coexist:
a \emph{physical} plane (\texttt{172.20.10.0/28}) carried by the wireless
interface, and a \emph{mesh} plane (\texttt{10.0.0.0/24}) carried by a TUN
virtual interface, through which all application traffic is routed by the
TDMA framework. The TDMA schedule uses a frame of $T_{\mathrm{frame}} =
150$\,ms divided into three slots of $T_{\mathrm{slot}} = 50$\,ms, one per
node, with a 5\,ms guard interval at the end of each slot. Slot~$i$ is
assigned to node~$i$, giving the fixed ordering
$[\text{N1}, \text{N2}, \text{N3}]$ shown in Figure~\ref{fig:frame}.

\begin{figure}[h]
\centering
\fbox{\parbox{0.9\textwidth}{\centering\vspace{1.2cm}
N1 slot $[0,50)$\,ms \quad|\quad N2 slot $[50,100)$\,ms \quad|\quad N3 slot $[100,150)$\,ms
\vspace{1.2cm}}}
\caption{TDMA frame structure with three synchronised nodes. The 150\,ms
frame is divided into three 50\,ms slots; each node transmits only within
its own slot.}
\label{fig:frame}
\end{figure}

All three nodes were placed on a laboratory bench within mutual radio range,
so that any pair could communicate directly. This placement is a deliberate
constraint of the testbed: because the nodes share a single collision domain,
link failures cannot be produced by distance and must instead be induced in
software (Section~\ref{subsec:linkfail}). The same limitation was reported by
Morais~\cite{morais2024synchronization}, and its consequences are discussed
where relevant.

\subsection{Hardware Platform}
\label{subsec:hardware}

Table~\ref{tab:hardware} lists the hardware of each node.

\begin{table}[h]
\centering
\caption{Hardware components of each testbed node}
\label{tab:hardware}
\begin{tabular}{llll}
\toprule
\textbf{Component} & \textbf{N1 (Robot)} & \textbf{N2 (Relay)} & \textbf{N3 (Base Station)} \\
\midrule
Platform        & AlphaBot2-Pi    & Raspberry Pi 4B & Laptop (Ubuntu 22.04) \\
Processor       & RPi 4B, 4\,GB   & RPi 4B, 4\,GB   & Intel Core i7 \\
Wireless        & Wi-Fi (ad-hoc)  & Wi-Fi (ad-hoc)  & Wi-Fi (ad-hoc) \\
Camera          & Pi Camera v2    & ---             & --- \\
Motor driver    & TB6612FNG       & ---             & --- \\
Servo controller & PCA9685 (I\textsuperscript{2}C) & --- & --- \\
Gamepad         & ---             & ---             & DualShock 4 (USB) \\
\bottomrule
\end{tabular}
\end{table}

\textbf{Node~1 (N1) --- Robot.} An AlphaBot2-Pi differential-drive robot with
a Raspberry Pi~4B and a Pi Camera Module~v2 mounted at the front. N1 runs the
RA-TDMAs+ node software and the \texttt{alphabot\_node.py} teleoperation
server, which manages video streaming, motor control, and telemetry.

\textbf{Node~2 (N2) --- Relay.} A Raspberry Pi~4B running only the RA-TDMAs+
node software. It originates and consumes no application data; it exists to
forward traffic between N1 and N3 when the direct link is unavailable.

\textbf{Node~3 (N3) --- Base Station.} A laptop running Ubuntu~22.04, running
the RA-TDMAs+ node software and \texttt{base\_station.py}, which displays the
incoming video and translates DualShock~4 gamepad input into teleoperation
commands. All measurements were taken at N3.

The AlphaBot2-Pi platform was chosen because it matches the hardware used by
Morais~\cite{morais2024synchronization}, enabling comparison of the two
routing approaches under identical physical conditions.

\subsection{Network and Software Configuration}
\label{subsec:netconfig}

Table~\ref{tab:netconfig} lists the address assignments. Each node sets its
wireless interface to ad-hoc mode on channel~6, assigns a physical address,
brings up the TUN interface with a mesh address, and enables IP forwarding.

\begin{table}[h]
\centering
\caption{Network address assignments}
\label{tab:netconfig}
\begin{tabular}{lccc}
\toprule
\textbf{Interface} & \textbf{N1} & \textbf{N2} & \textbf{N3} \\
\midrule
\texttt{wlan0} (physical) & \texttt{172.20.10.1} & \texttt{172.20.10.2} & \texttt{172.20.10.3} \\
\texttt{tun} (mesh)       & \texttt{10.0.0.1}    & \texttt{10.0.0.2}    & \texttt{10.0.0.3}    \\
\bottomrule
\end{tabular}
\end{table}

Two distinct traffic planes travel over the air, and they must not be
confused when interpreting the captures:

\begin{itemize}
  \item \textbf{Synchronisation beacons (control).} Each node broadcasts one
  beacon per frame, \emph{within its own slot}, to its neighbours as a small
  UDP packet (physical addresses, ports 7001--7003, $\approx 93$--$107$\,B).
  These beacons carry the adjacency matrix and slot-timing information used by
  the RA-TDMAs+ synchroniser. They are \emph{not} relayed: every node emits
  its own.
  \item \textbf{Application data.} The video stream and the control channel
  travel over the mesh plane (\texttt{10.0.0.x}). The video is an MPEG-TS/UDP
  flow to \texttt{10.0.0.3:5000}; the teleoperation RTT probe is a UDP echo to
  \texttt{10.0.0.1:19876}. When the direct link is unavailable, this traffic
  is forwarded through N2.
\end{itemize}

\subsection{Video Pipeline and Application Parameters}
\label{subsec:video-pipeline}

On N1, \texttt{rpicam-vid} captures raw YUV~4:2:0 frames from the camera and
pipes them to \texttt{ffmpeg}, which encodes them with \texttt{libx264},
wraps the output in MPEG-TS, and sends each transport-stream segment as a UDP
datagram to \texttt{10.0.0.3:5000}. Because the destination is a mesh (TUN)
address, the datagram enters the TDMA routing layer transparently. On N3,
\texttt{base\_station.py} receives each datagram, timestamps it for
measurement, and forwards it to \texttt{ffplay}. Table~\ref{tab:app-params}
lists the parameters.

\begin{table}[h]
\centering
\caption{Application parameters used in all experiments}
\label{tab:app-params}
\begin{tabular}{lll}
\toprule
\textbf{Parameter} & \textbf{Value} & \textbf{Source} \\
\midrule
Resolution           & $640 \times 480$\,px  & \texttt{alphabot\_node.py} \\
Frame rate           & 15\,fps               & \texttt{VIDEO\_FPS} \\
Target bitrate       & 500\,kbps             & \texttt{VIDEO\_BITRATE} \\
Encoder / preset     & \texttt{libx264} / \texttt{ultrafast}, \texttt{zerolatency} & \\
GOP size             & 1 (all-intra)         & \texttt{-g 1} \\
Container            & MPEG-TS               & \texttt{-f mpegts} \\
MPEG-TS UDP size     & 1316\,B ($7\times188$)& \texttt{pkt\_size=1316} \\
Video destination    & \texttt{10.0.0.3:5000}& \texttt{BASE\_IP:VIDEO\_PORT} \\
Command port / rate  & UDP 9000 / 20\,Hz     & \texttt{CMD\_PORT} \\
Telemetry port / rate& UDP 9001 / 5\,Hz      & \texttt{TEL\_PORT} \\
RTT echo port        & UDP 19876             & \texttt{rtt\_test.py} \\
\bottomrule
\end{tabular}
\end{table}

\subsection{Link-Failure Methodology}
\label{subsec:linkfail}

Because all nodes share a collision domain, a direct-link failure is induced
logically rather than physically. At the moment of failure, the following
rules are applied on N3:

\begin{verbatim}
iptables -I INPUT  1 -s 172.20.10.1 -j DROP
iptables -I OUTPUT 1 -d 172.20.10.1 -j DROP
\end{verbatim}

These rules discard, at the IP layer, every frame to and from N1's physical
address. This makes the direct link invisible to N3's TDMA framework and to
the TCP data plane, forcing traffic onto the relay path N1\,$\to$\,N2\,$\to$\,N3.
The rules are cleared between repetitions with \texttt{iptables -F}.

A subtlety of this method, relevant to the synchronisation analysis of
Section~\ref{subsec:res-slots}, is that \texttt{iptables} filters at the IP
layer, whereas a packet \emph{capture} (\texttt{tcpdump}) taps the interface
\emph{before} the firewall. Consequently, a capture taken at N3 with the link
blocked still records N1's beacons (they reach the antenna), even though N3's
TDMA application no longer receives them. The block therefore changes which
node N3 synchronises to (it falls back to synchronising through N2), without
removing N1's transmissions from the medium. To verify that the block was
effective for the data plane, each relay measurement was preceded by a check
that N1's physical address was unreachable (\texttt{ping} to
\texttt{172.20.10.1} failing), since the physical address, unlike the mesh
address, is not reachable through the relay.


%% =======================================================================
\section{Metrics}
\label{sec:metrics}

\paragraph{Round-trip time (RTT).}
$\mathrm{RTT}_k = t^{\mathrm{rx}}_k - t^{\mathrm{tx}}_k$, the elapsed time from
N3 sending a UDP echo request to receiving N1's reply. In a TDMA network the
RTT reflects the slot structure: each leg of the round trip can depart only
within its node's slot.

\paragraph{Packet delivery ratio (PDR).}
$\mathrm{PDR} = N_{\mathrm{rx}}/N_{\mathrm{tx}} \times 100\,\%$, the fraction
of echo requests answered within a timeout.

\paragraph{Jitter.}
$J = \frac{1}{N-1}\sum_{k=2}^{N}\lvert \mathrm{RTT}_k - \mathrm{RTT}_{k-1}\rvert$,
the mean absolute deviation of consecutive RTT samples.

\paragraph{Throughput.}
$T = 8\,B_{\mathrm{total}} / \Delta t$ [kbps], where $B_{\mathrm{total}}$ is
the sum of UDP video payload bytes received in an observation window
$\Delta t$.

\paragraph{Convergence time.}
$t_{\mathrm{conv}} = t_{\mathrm{first\,relay}} - t_{\mathrm{last\,direct}}$,
the gap between the last video packet on the direct path and the first packet
recovered through the relay.

\paragraph{Inter-arrival time (IA).}
$\mathrm{IA}_k = t_{k+1} - t_k$, the time between consecutive video packets.
Its distribution reveals the TDMA burst structure and allows verification
that relay forwarding preserves it.

\paragraph{Slot occupancy.}
For the synchronisation analysis, the transmission instant of each captured
beacon is wrapped to the frame, $\phi = t \bmod T_{\mathrm{frame}}$, and
attributed to its source node. A node \emph{respects its slot} if its
beacons cluster within its assigned 50\,ms interval. Because the RA-TDMAs+
synchroniser advances all slots by a common amount each round (it does not
compensate the transmission/processing latency), the whole frame rotates
slowly in wall-clock time; this common rotation is removed before evaluating
the per-node clustering, leaving the \emph{relative} synchronisation between
nodes.


%% =======================================================================
\section{Experiments}
\label{sec:experiments}

\textbf{Experiment 1 --- Control-channel RTT and PDR.} N3 sends 100 UDP echo
requests to \texttt{10.0.0.1} at 10\,Hz; N1 echoes each immediately. The 100
samples are recorded at N3. The experiment is run under direct-link and relay
conditions. It characterises the control-plane latency and verifies that it
follows the TDMA slot structure.

\textbf{Experiment 2 --- Video throughput.} N1 streams video continuously to
N3. A passive \texttt{tcpdump}/raw-socket capture at N3 (which does not occupy
port~5000 and therefore does not perturb \texttt{ffplay}) records every video
packet for 15\,s. Throughput is computed over the window. The experiment is
run under direct-link and relay conditions, with three repetitions each.

\textbf{Experiment 3 --- Routing convergence.} Starting from a stable
direct-link stream, the \texttt{iptables} rules of
Section~\ref{subsec:linkfail} are applied; the capture is used to detect the
video gap and the first packet recovered via N2, giving $t_{\mathrm{conv}}$.
Three repetitions per run, two runs.

\textbf{Experiment 4 --- Packet inter-arrival timing.} The inter-arrival
times of video packets are extracted from the captures of Experiment~2,
separately for direct and relay, over 30\,s. Histograms are produced and
compared to verify that the relay preserves the TDMA burst structure.

\textbf{Experiment 5 --- TDMA slot synchronisation.} The beacons emitted by
the three nodes are captured and their transmission instants are wrapped to
the frame to verify that each node occupies its assigned slot. The analysis
is performed under direct conditions and repeated under relay conditions to
confirm that synchronisation is maintained while the relay forwards the video.


%% =======================================================================
\section{Results and Discussion}
\label{sec:results}

\subsection{Control-Channel Latency}
\label{subsec:res-rtt}

Figures~\ref{fig:rtt-direct} and~\ref{fig:rtt-relay} show the RTT of the
control channel over the direct link and the relay path, as a per-ping
sequence (left) and as a distribution (right). In both cases the PDR is
100\,\% and no single mean or median is representative, because the
distribution is strongly \emph{bimodal}; Table~\ref{tab:rtt} therefore
characterises each cluster separately.

\begin{figure}[H]
\centering
\includegraphics[width=\textwidth]{RTT_over_the_direct_link.png}
\caption{RTT of the control channel over the direct link. Left: RTT per ping
sequence; the values alternate in blocks between a low cluster
($\approx 50$\,ms) and a high cluster ($\approx 200$\,ms). Right: the
resulting bimodal distribution, with the two clusters separated by one TDMA
frame (150\,ms).}
\label{fig:rtt-direct}
\end{figure}

\begin{figure}[H]
\centering
\includegraphics[width=\textwidth]{RTT_over_the_relay_path.png}
\caption{RTT of the control channel over the relay path
(N3\,$\to$\,N2\,$\to$\,N1). The distribution is practically identical to the
direct case, showing that relaying adds no significant latency to the control
channel.}
\label{fig:rtt-relay}
\end{figure}

\begin{table}[h]
\centering
\caption{Control-channel RTT (N3 $\to$ N1), bimodal characterisation
($n = 100$ per scenario, PDR $= 100\,\%$)}
\label{tab:rtt}
\begin{tabular}{lcc}
\toprule
\textbf{Metric} & \textbf{Direct link} & \textbf{Relay path} \\
\midrule
Lower cluster (catches the round)   & $\approx 52$\,ms  & $\approx 52$\,ms  \\
Upper cluster (misses by one round) & $\approx 202$\,ms & $\approx 201$\,ms \\
Cluster separation                  & $\approx 150$\,ms & $\approx 150$\,ms \\
Minimum RTT                         & 42.6\,ms          & 36.5\,ms          \\
Maximum RTT                         & 250\,ms           & 229\,ms           \\
Jitter                              & 15.2\,ms          & 14.6\,ms          \\
PDR                                 & 100\,\%           & 100\,\%           \\
\bottomrule
\end{tabular}
\end{table}

The RTT does not concentrate around a single value; it splits into two
well-separated clusters with an empty band between roughly 60 and 190\,ms.
This is a direct consequence of the TDMA slot structure. A round trip
traverses two slot waits: the request waits for N3's slot, and the reply
waits for N1's slot. Because each leg can depart only within its node's slot,
the reply either reaches N1's transmit opportunity in the immediately
following frame --- giving the lower cluster ($\approx 52$\,ms, set by the
fixed encapsulation and the once-per-frame transmission overhead) --- or
narrowly misses it and must wait one entire frame, giving the upper cluster
($\approx 52 + 150 = 202$\,ms). The cluster separation of exactly one frame
(150\,ms) is the signature of TDMA round quantisation: latency is lost only in
whole-frame units, which is why no values fall between the clusters. The
proportion of pings in each cluster varies between runs, because the
free-running 10\,Hz probe cadence beats against the 150\,ms frame; for this
reason the mean and median are not reported as headline figures.

Crucially, the 100\,\% PDR and the bounded maximum (below two frames) confirm
that the slots are \emph{respected}. A desynchronised network would instead
exhibit unbounded, unstructured latency and packet loss; the clean two-level
quantisation observed here is positive evidence that the schedule holds.

Normalising by the frame period, the mean RTT of $\approx 0.9$ frame is
consistent with the value reported by Morais~\cite{morais2024synchronization}
(88\,ms over a 100\,ms frame, i.e.\ 0.88 frame), confirming that both systems
exhibit the same TDMA-imposed latency behaviour. It is worth noting that an
earlier configuration of the system, in which the transmit slot was derived
from the absolute wall-clock rather than from the beacon-adjusted slot,
produced an apparent mean RTT of $\approx 50$\,ms (0.33 frame). That value is
\emph{below} the reference mean and is not representative: it corresponded to
the slot being locked, by chance, in a favourable phase in which the reply
always caught the following slot. Closing the synchronisation loop exposed the
true, phase-independent distribution reported here.

The relay path produces a distribution practically identical to the direct
path (Table~\ref{tab:rtt}, Figure~\ref{fig:rtt-relay}), with the same two
clusters and the same separation. The relay therefore adds no significant
latency to the control channel, in agreement with the throughput result of
the next section: kernel IP forwarding at N2 forwards the packets within the
same slot in which it receives them.

\subsection{Video Throughput}
\label{subsec:res-throughput}

Table~\ref{tab:throughput} presents the video throughput measured at N3 under
direct-link and relay conditions, averaged over three repetitions.

\begin{table}[h]
\centering
\caption{Video stream throughput at N3 (passive capture, 15\,s window)}
\label{tab:throughput}
\begin{tabular}{lccc}
\toprule
\textbf{Condition} & \textbf{Throughput} & \textbf{Mean packet size} & \textbf{Packet rate} \\
\midrule
Direct & 558\,kbps & 1135\,B & 61\,pkt/s \\
Relay  & 552\,kbps & 1114\,B & 61\,pkt/s \\
\midrule
\textbf{Overhead (relay vs.\ direct)} & \multicolumn{3}{c}{$\approx 1\,\%$} \\
\bottomrule
\end{tabular}
\end{table}

The direct-link mean of 558\,kbps slightly exceeds the 500\,kbps encoder
target for two reasons. First, the MPEG-TS container adds $\approx 15\,\%$
overhead (a 4-byte header per 188-byte transport packet). Second, the encoder
is configured with \texttt{-g~1}, disabling inter-frame prediction; every
frame is an independent intra frame, producing a higher and more constant
bitrate than a configuration with P/B-frames.

Under relay conditions the mean throughput is 552\,kbps, a reduction of only
$\approx 1\,\%$ relative to the direct case. This difference is within the
run-to-run variability of the H.264 encoder as scene content changes; no
systematic degradation attributable to the relay hop is observed. The relay
is therefore effectively \emph{transparent} for the video plane: the kernel
\texttt{ip\_forward} path at N2 carries the stream without consuming a
separate TDMA scheduling step. This is the data-plane counterpart of the RTT
result, where the relay likewise added no measurable latency.

\subsection{Routing Convergence}
\label{subsec:res-convergence}

Table~\ref{tab:convergence} presents the convergence times across the six
link-failure repetitions of Experiment~3.

\begin{table}[h]
\centering
\caption{Routing convergence time after a simulated direct-link failure}
\label{tab:convergence}
\begin{tabular}{ccl}
\toprule
\textbf{Run} & \textbf{$t_{\mathrm{conv}}$ (ms)} & \textbf{Notes} \\
\midrule
1 & 2675 & --- \\
2 & 2863 & --- \\
3 & 2926 & --- \\
4 & 2933 & --- \\
5 & 453  & outlier (relay route still warm; near-instant failover) \\
6 & ---  & no gap detected (warm route; failover $< 0.4$\,s threshold) \\
\midrule
\multicolumn{2}{l}{\textbf{Median (all valid)}} & \textbf{2895\,ms} \\
\bottomrule
\end{tabular}
\end{table}

The median convergence time is $\approx 2.9$\,s, equivalent to about 19~TDMA
frames. The value is highly consistent across the well-behaved repetitions
(2675--2933\,ms). Two repetitions are atypical and are reported separately:
in both, the relay route was still installed from a previous repetition, so
the failover completed in less than the 0.4\,s detection threshold and the
video stream barely interrupted (one of these registered 453\,ms, the other
showed no measurable gap). The median is used as the headline value because
it is robust to these outliers.

The recovery chain, after the direct link is severed, comprises:

\begin{enumerate}
  \item \textbf{Beacon-loss detection ($\approx 450$\,ms).} RA-TDMAs+ declares
  a link failure after three consecutive missed beacons; at a 150\,ms frame
  this is $3 \times 150 = 450$\,ms.
  \item \textbf{Route recomputation and installation ($\approx 100\,\mu$s).}
  The link quality to N1 drops to zero, the spanning tree is recomputed, and
  the host routes are reinstalled in the kernel via Netlink --- all in
  microseconds on the three-node graph.
  \item \textbf{TCP teardown and reconnection via the relay ($\approx 2.4$\,s).}
  The data-plane TCP connection established over the direct path must be torn
  down (its retransmission timeout must expire) and re-established over the
  relay path before video resumes.
\end{enumerate}

The dominant contribution is the TCP teardown and reconnection, not the
routing computation. This is an important architectural observation: in the
Layer~3 design the transport connection is bound to a path, so a path change
forces a transport reset. The implication for the comparison of
Section~\ref{sec:comparison} is direct, and reducing the TCP retransmission
and keep-alive timers is the most effective way to shorten convergence.

\subsection{Packet Inter-Arrival Distribution}
\label{subsec:res-ia}

Figures~\ref{fig:ia-direct} and~\ref{fig:ia-relay} show the inter-arrival
histograms of video packets at N3 under direct and relay conditions, and
Table~\ref{tab:ia} gives the statistics.

\begin{figure}[H]
\centering
\includegraphics[width=0.8\textwidth]{ia_direct.png}
\caption{Inter-arrival histogram of video packets at N3, direct link. The
distribution is strongly bimodal: a peak at $<10$\,ms (packets within one
TDMA burst) and a peak at $\approx 140$--$150$\,ms (the inter-burst gap of one
frame).}
\label{fig:ia-direct}
\end{figure}

\begin{figure}[H]
\centering
\includegraphics[width=0.8\textwidth]{ia_relay.png}
\caption{Inter-arrival histogram of video packets at N3, relay path. The
bimodal structure is preserved identically, confirming that N2's kernel IP
forwarding introduces no inter-packet delay at the burst level.}
\label{fig:ia-relay}
\end{figure}

\begin{table}[h]
\centering
\caption{Video inter-arrival statistics at N3 (30\,s window)}
\label{tab:ia}
\begin{tabular}{lcc}
\toprule
\textbf{Metric} & \textbf{Direct} & \textbf{Relay} \\
\midrule
Mean IA (ms)    & 16.3   & 16.3   \\
Median IA (ms)  & 0.7    & 0.7    \\
Maximum IA (ms) & 196    & 300    \\
P95 IA (ms)     & 145    & 145    \\
\% IA $< 10$\,ms & 88\,\% & 88\,\% \\
\bottomrule
\end{tabular}
\end{table}

The distribution is strongly bimodal in both conditions: $\approx 88\,\%$ of
inter-arrival intervals fall below 10\,ms (packets within the same burst) and
$\approx 11\,\%$ fall in the 130--150\,ms range (the inter-burst gap, one
frame, during which N2 and N3 hold the channel). This arises directly from
the schedule: N1 accumulates encoded data for $\approx 100$\,ms and then
transmits a burst during its 50\,ms slot. At 558\,kbps with 1316-byte MPEG-TS
UDP packets, each frame's worth of data is
\begin{equation}
n_{\mathrm{pkts}} = \frac{558{,}000 \times 0.150}{8 \times 1316} \approx 8\,\text{packets},
\label{eq:pkts-per-slot}
\end{equation}
consistent with the observed burst size of 8--9 packets.

The relay histogram is statistically indistinguishable from the direct one:
the median, P95, and burst fraction are identical, and only the rare tail
differs. The relay shifts the absolute arrival time of each burst by about
one slot but leaves the intra-burst and inter-burst timing unchanged. This is
the clearest demonstration that the relay forwards each packet within the
same slot in which it arrives, preserving the TDMA structure of the stream.
It is the streaming-data counterpart of the beacon analysis below.

\subsection{TDMA Slot Synchronisation}
\label{subsec:res-slots}

The previous sections show that the \emph{end-to-end timing} is consistent
with the schedule. This section proves the schedule directly, by examining
the instant at which each node transmits. Whereas the video stream is
generated only by N1 (and therefore only evidences N1's slot), the
synchronisation beacons are emitted by all three nodes and thus reveal the
complete slot structure. This is the same method used by
Morais~\cite{morais2024synchronization}.

Figure~\ref{fig:slots} wraps the beacon transmission instants to the 150\,ms
frame over a sequence of rounds. The three nodes occupy three disjoint
regions in the strict cyclic order N1\,$\to$\,N2\,$\to$\,N3, each separated
from the next by approximately one slot (50\,ms), and each node stays in its
region across rounds. Quantitatively, after removing the common frame
rotation, the per-node standard deviation of the transmission instant is
$\approx 16$\,ms and 84--89\,\% of beacons fall within the assigned 50\,ms
slot.

\begin{figure}[H]
\centering
\includegraphics[width=0.85\textwidth]{slot_occupancy.png}
\caption{Beacon transmission instant wrapped to the TDMA frame, per round.
Each node occupies its own slot in the cyclic order N1\,$\to$\,N2\,$\to$\,N3,
stable across rounds --- direct evidence that the nodes respect the schedule.}
\label{fig:slots}
\end{figure}

Two characteristics of the RA-TDMAs+ synchroniser must be noted when reading
this figure. First, the absolute phase of the frame is arbitrary: there is no
global clock, so only the \emph{relative} ordering and spacing of the nodes
carry meaning. Second, the synchroniser does not compensate the
transmission/processing latency, so all slots advance by a small common
amount each round and the whole frame rotates slowly in wall-clock time (here
$\approx 18$\,ms/s). This rotation is a property of the algorithm, also
present in the reference work, and is removed before evaluating the per-node
clustering; over a short window of a few rounds the rotation is negligible and
the slots appear stationary, as in Figure~\ref{fig:slots}.

Under relay conditions, the same beacon analysis confirms that the cyclic slot
ordering is maintained. A subtlety specific to the link-failure method
(Section~\ref{subsec:linkfail}) is relevant here: with the direct link
blocked at N3, N3's TDMA application no longer \emph{receives} N1's beacons
(they are discarded by \texttt{iptables} before the socket), so N3
synchronises through the two-hop chain N1\,$\to$\,N2\,$\to$\,N3. The packet
capture, however, still records N1's beacons, because \texttt{tcpdump} taps
the interface before the firewall and all nodes share the medium. The slots
therefore remain occupied and ordered; the only effect of the relay is a
slightly looser alignment of N3 relative to N1, consistent with N3
synchronising indirectly through N2 rather than directly.

Together, the inter-arrival analysis (streaming data, N1's slot) and the
beacon analysis (all three slots) establish the central claim of this
evaluation: the nodes respect their TDMA slots, for both periodic control
traffic and continuous video, and this property is preserved under relaying.


%% =======================================================================
\section{Comparison with the Layer~2 Approach}
\label{sec:comparison}

The system evaluated here implements routing at Layer~3 (IP), using kernel
forwarding and standard routing tables over a TUN overlay. The prior work of
Morais~\cite{morais2024synchronization} implements routing at Layer~2, by
manipulating ARP tables so that the relay forwards frames transparently below
IP. Both run on the same RA-TDMAs+ synchroniser and the same class of
hardware. Table~\ref{tab:comparison} summarises the comparison; the
quantitative latency figures are normalised to the frame period to account
for the different frame durations (100\,ms with four nodes in the reference,
150\,ms with three nodes here).

\begin{table}[h]
\centering
\caption{Comparison between the Layer~2 (reference) and Layer~3 (this work)
routing approaches}
\label{tab:comparison}
\begin{tabular}{lll}
\toprule
\textbf{Aspect} & \textbf{Layer~2 (Morais)} & \textbf{Layer~3 (this work)} \\
\midrule
Forwarding mechanism & ARP-table manipulation & IP forwarding (TUN + routes) \\
Routing layer        & Data link             & Network \\
Mean RTT (norm.)     & 0.88 frame            & $\approx 0.9$ frame \\
RTT distribution     & spread / triangular   & bimodal (frame-quantised) \\
PDR                  & 100\,\%               & 100\,\% \\
Slot compliance      & verified (beacons)    & verified (beacons + video) \\
Path observability   & invisible at IP layer & visible (\texttt{ip route}, \texttt{traceroute}) \\
Convergence          & faster (no transport reset) & $\approx 2.9$\,s (TCP teardown dominates) \\
Reported stability   & runtime failures      & stable across all experiments \\
\bottomrule
\end{tabular}
\end{table}

The two approaches achieve equivalent TDMA-level behaviour: the normalised RTT
(0.9 vs.\ 0.88 frame), the 100\,\% PDR, and the verified slot compliance are
essentially the same, as expected since both share the synchroniser and the
slot structure. The differences are architectural.

\textbf{Latency distribution.} The reference reports a spread, triangular RTT
distribution, whereas this work observes a clean bimodal distribution. Both
are manifestations of the same slot-waiting mechanism; the bimodal shape
arises here because, with three nodes and the requester/responder slots
adjacent across the frame boundary, the reply either catches the next slot or
misses it by exactly one frame, producing two discrete levels rather than a
continuum.

\textbf{Convergence.} This is the principal disadvantage of the Layer~3
approach. Because the data-plane TCP connection is bound to a path, a link
failure forces a transport teardown and reconnection, which dominates the
$\approx 2.9$\,s recovery time. The Layer~2 approach changes the path below
IP, so the transport connection survives a route change transparently and
recovers within a frame or two. Mitigating this gap --- for instance by
shortening the TCP timers or by using a connectionless transport for the
relayed data --- is the clearest avenue for improving the Layer~3 design.

\textbf{Observability and stability.} The Layer~3 approach trades convergence
speed for two practical advantages. The relay path is visible to standard IP
tools, which simplifies debugging and operation, whereas the Layer~2 path is
invisible at the IP layer. And, in contrast to the runtime instability
reported for the Layer~2 implementation, the Layer~3 system completed every
experiment without failure, benefiting from relying on mature kernel
forwarding rather than on ARP manipulation.


%% =======================================================================
\section{Summary}
\label{sec:eval-summary}

This chapter evaluated the Layer~3 routing system across five experiments on
physical hardware. The control channel exhibits a bimodal RTT whose two
clusters are separated by exactly one TDMA frame and whose mean of
$\approx 0.9$ frame matches the reference work, with 100\,\% PDR; the
quantisation and the absence of loss are positive evidence that the slots are
respected. Video throughput under relaying is within $\approx 1\,\%$ of the
direct baseline, and the inter-arrival distribution is identical in both
conditions, confirming that kernel IP forwarding at the relay preserves the
TDMA burst structure of the stream. The beacon analysis shows the three nodes
occupying their slots in strict cyclic order, both directly and through the
relay, establishing the central claim that the nodes respect their TDMA slots
for periodic control traffic and for continuous video alike. Routing
convergence after a link failure has a median of $\approx 2.9$\,s, dominated
by TCP teardown and reconnection rather than by routing computation. The
comparison with the Layer~2 approach of Morais~\cite{morais2024synchronization}
shows equivalent TDMA-level performance, a faster convergence for Layer~2 due
to its transparent path change, and better observability and stability for the
Layer~3 design developed here. Reducing the convergence time, and completing a
side-by-side re-run of the Layer~2 system on the same testbed, are identified
as the main directions for future work.
